\chapter{Hypothesis Testing}

% -------------------------------------------------------
\section{Hypotheses}
% -------------------------------------------------------

\begin{definition}[Null Hypothesis]
The \textbf{null hypothesis} $H_0$ is the default assumption — typically a statement of
no effect, no difference, or a specific parameter value. It is the hypothesis we attempt
to reject.
\end{definition}

\begin{definition}[Alternative Hypothesis]
The \textbf{alternative hypothesis} $H_1$ (or $H_a$) is the claim we seek evidence for.
It is accepted when $H_0$ is rejected.
\end{definition}

% -------------------------------------------------------
\section{Types of Error}
% -------------------------------------------------------

\begin{table}[h!]
\centering
\renewcommand{\arraystretch}{1.5}
\begin{tabular}{lcc}
\toprule
 & \textbf{$H_0$ is True} & \textbf{$H_0$ is False} \\
\midrule
\textbf{Reject $H_0$}     & Type I error ($\alpha$) & Correct decision \\
\textbf{Fail to reject $H_0$} & Correct decision    & Type II error ($\beta$) \\
\bottomrule
\end{tabular}
\caption{Decision outcomes in hypothesis testing.}
\end{table}

\begin{definition}[Type I Error]
Rejecting $H_0$ when it is in fact \emph{true}. The probability of a Type I error is the
\textbf{significance level} $\alpha$.
\end{definition}

\begin{definition}[Type II Error]
Failing to reject $H_0$ when it is in fact \emph{false}. The probability of a Type II
error is $\beta$.
\end{definition}

% -------------------------------------------------------
\section{Significance Level}
% -------------------------------------------------------

\begin{definition}[Significance Level]
The \textbf{significance level} $\alpha$ is the maximum acceptable probability of
committing a Type I error. A common choice is
\[
  \alpha = 0.05 \quad (5\%).
\]
\end{definition}

% -------------------------------------------------------
\section{Test Statistics}
% -------------------------------------------------------

\subsection{Z-Test}

Used when the population standard deviation $\sigma$ is \emph{known}, or when the
sample size $n$ is large ($n \geq 30$).

\begin{definition}[Z-Test Statistic]
\[
  Z = \frac{\bar{X} - \mu}{\sigma / \sqrt{n}},
\]
where $\bar{X}$ is the sample mean, $\mu$ is the hypothesised population mean,
$\sigma$ is the population standard deviation, and $n$ is the sample size.
Under $H_0$, $Z \sim \mathcal{N}(0,1)$.
\end{definition}

\subsection{$t$-Test}

Used when $\sigma$ is \emph{unknown} and the sample size is small.

\begin{definition}[$t$-Test Statistic]
\[
  T = \frac{\bar{X} - \mu}{s / \sqrt{n}},
\]
where $s$ is the sample standard deviation. Under $H_0$, $T \sim t_{n-1}$
(Student's $t$-distribution with $n-1$ degrees of freedom).
\end{definition}

% -------------------------------------------------------
\section{Decision Rule}
% -------------------------------------------------------

\begin{theorem}[Decision Rule]
There are two equivalent formulations:
\begin{enumerate}
  \item \textbf{Critical value approach}: Reject $H_0$ if the computed test statistic
        exceeds the critical value at significance level $\alpha$.
  \item \textbf{$p$-value approach}: Reject $H_0$ if the $p$-value $\leq \alpha$.
\end{enumerate}
\end{theorem}

\begin{remark}
The $p$-value is the probability of observing a test statistic at least as extreme as
the one computed, assuming $H_0$ is true. A small $p$-value provides strong evidence
against $H_0$.
\end{remark}
